\section{Introduction}
\label{sec:intro}

% Target length: ~1 page. Keep paragraphs tight; the course rubric
% explicitly penalizes verbose, low-information writing.

\paragraph{Motivation.}
% 1 paragraph: ViTs have weaker inductive biases than CNNs (no locality,
% no translation equivariance), so they typically need more data or
% stronger regularization to match CNN generalization
% \citep{dosovitskiy2021vit,touvron2021deit}. On small datasets the
% problem is sharper: the model can memorize the training set long before
% it learns the right features.
\draftnote{Write 4--6 sentences motivating regularization for ViTs.}

\paragraph{Why CIFAR-100 with ViT-Tiny.}
% 1 paragraph: CIFAR-100 (50k train, 100 classes) is small relative to
% ImageNet but large enough that fine-tuning a small ViT is non-trivial.
% ViT-Tiny is the smallest standard ViT and therefore most sensitive to
% regularizer strength: too weak and it overfits; too strong and it
% underfits within 30 epochs. This regime is informative.
\draftnote{One paragraph: justify the dataset/model choice as a
realistic ``modest model, modest data'' overfitting regime.}

\paragraph{Contributions.}
This paper makes the following contributions:
\begin{itemize}
  \item We establish a unified, reproducible fine-tuning baseline for
        ViT-Tiny/16 on CIFAR-100 against which all regularizers are
        compared (\cref{sec:baseline}).
  \item We perform a systematic sweep of \emph{weight decay} strengths
        in AdamW and analyze its effect on the train--validation gap
        (\cref{sec:wd}).
  \item We study \emph{dropout} placement and rate (\cref{sec:dropout}).
  \item We compare a suite of \emph{data augmentations} from light to
        heavy combinations, including Mixup/CutMix
        (\cref{sec:augmentation}).
  \item We study \emph{label smoothing} and \emph{early stopping} as
        output-level and training-process regularizers
        (\cref{sec:lsearly}).
  \item We provide a unified cross-technique comparison and discuss
        which family of regularization matters most for ViT in this
        regime (\cref{sec:comparison,sec:conclusion}).
\end{itemize}

\paragraph{Roadmap.}
\Cref{sec:related} reviews related work.
\Cref{sec:methodology} defines the baseline and the regularization
families under study.
\Cref{sec:experiments} describes the shared experimental protocol and
reports results for each method.
\Cref{sec:comparison} compares them.
\Cref{sec:conclusion} concludes.
